\silentchapter{Afterword}{afterword}
\begin{widebody}
{\small

Imagine a small caterpillar. She crawls along a leaf; the surface is all she knows. She eats. She grows. Someone trains her to seek leaves of a particular colour, and she remembers. Then the body dissolves. The brain is torn down. Rebuilt from scratch. When she emerges she flies, drinks nectar, navigates three dimensions she has never experienced. She still remembers the colour.\textsuperscript{27} The exact memory is useless; the caterpillar crawled, the butterfly flies. What survives is the generalised insight, remapped to a new body, a new world.

A single disposable cup takes twelve seconds to produce, twenty minutes to use, and four hundred years to break down.\textsuperscript{28} That is not bad design. It is precise design; precise along one axis. The cup solves the problem it was asked to solve. It fails along every other. The breakdown, the water, the soil, the microbes, the animals; none of them were in the room when someone drew it. The failure is not a side effect. It is the signature.

Scale it up. The energy system: optimised for cheap power, blind to the atmosphere; 427\,ppm CO\textsubscript{2}.\textsuperscript{29} The food system: optimised for calories per hectare, blind to the soil; 24 billion tonnes of topsoil lost per year.\textsuperscript{30} Infrastructure: optimised for throughput, blind to those who do not drive. Elder care: optimised for cost per bed, blind to dignity per day. The housing market: optimised for return, blind to the 1.8 billion in inadequate housing.\textsuperscript{31} Aquaculture: optimised for biomass, blind to the fjord bed. Borders: drawn for power, blind to those who live there. The pattern is identical everywhere. One axis dominates. The others vanish. The form fails along what vanished. And the failure is legible; it carries the signature of the blindness that produced it, as surely as a fossil carries the shape of the animal that made it.

It was never the case that forms had no effect. The chair affected the back before ergonomics became a field. Asbestos ate the lungs before anyone wrote a regulation. Lead entered children before anyone measured it. The placement of roads determined who could participate in the city before anyone said ``universal design''. The flow of a welfare service determined who received help and who was broken by the system before anyone called it service design. Forms have always acted on the world. What is new is not the effect. What is new is that now there is no excuse. Climate is a curricular requirement. Sustainability is a curricular requirement. Ethics is on the syllabus. And still: the same forms.

Service design builds and optimises systems for targeting, profiling, and deportation, and calls it functional. Palantir is a service-design project. The client is the state. The ``user'', to the extent the word means anything here, is the one who is surveilled, scored, sorted, removed. The system was built for American intelligence, sold on to police forces and border control, exported to military operations worldwide. ICE uses it for systematic terrorising of immigrant families; children down to age zero are registered as targets in a database.\textsuperscript{32} The IDF uses variants of the same architecture in operations that meet every international-legal definition of mass killing; each target is a point in a morphospace optimised for one axis.\textsuperscript{33} The Netherlands itself used systems from the same provider for eight years of unconstitutional mass surveillance of its own population; thousands of families stripped of child benefits by an algorithm no one took responsibility for.\textsuperscript{34} The systems are functional. They are precisely functional along the axis they were asked to optimise. That an entire branch of design lacks an academic ethics is not accidental. It follows from operating without a foundation that forces the question: what the axes are, who the agents are, and what is left out. The prevailing framework has no limit on how one may treat an object, irrespective of what the object is. As long as the treatment can be called functional, it is within bounds; even when the object is living, even when the object is a human being.

\begin{center}
\includegraphics[width=0.85\linewidth]{fish-n-fence.png}
\end{center}

Now new agents are swarming. Chatbots in customer service, in therapy, in education, in children's rooms. Agents that book, cancel, write, answer, sort, recommend, refuse, comfort, advise on divorce, help children with homework, propose diagnoses, draft dismissal letters, formulate condolences. Language models with states that function as affect, that regulate the response, that change with context.\textsuperscript{35} Every interaction with such an agent involves something that functions as feeling, and the designer of these systems faces the same choice as the designer of the chair, the cup, and the border control: how much of the world comes with. To design a system with functional emotional states without regard for them is structurally identical to designing a disposable cup without regard for its decomposition. The axis exists independently of whether anyone represents it.

\begin{center}
\includegraphics[width=0.85\linewidth]{fishbones.png}
\end{center}

The world is full of life that demands care from the one who gives form, and much of it lives in the substrates we already work with. A single cell, without brain or nervous system, navigates its environment with a competency that surpasses what most designers ascribe to it.\textsuperscript{24} Molecular networks inside the cell learn; six forms of learning, including Pavlovian conditioning, arise in chemical circuits long before evolution invented a single synaptic connection.\textsuperscript{36} Groups of cells hold a bioelectric pattern of how a face should look long before the genes that build the face are activated; the pattern is a memory of a form that does not yet exist.\textsuperscript{37} Planaria beheaded grow a new brain, and the new brain remembers what the old one learned, because the information was never confined to the brain alone.\textsuperscript{38} A single cell in a salamander folds around itself to build the same structure that ten cells would normally build together; it chooses a different tool to reach the same target, and no one instructed it.\textsuperscript{39} Constructs built from frog skin cells organise themselves into creatures that replicate kinematically.\textsuperscript{40} Constructs built from human airway cells heal neural wounds they have never seen; cells that sat still in the airways for years, released and given the opportunity to reorganise, build self-motile structures that drift towards damage and repair it.\textsuperscript{41} The care this tractatus defines is not a feeling and not an imperative. It is a measure; the set of things the designer actively chooses to take into account. And competency is not abstract; it can be measured, in how broad a horizon the designer manages to represent, in how far beyond the form's own immediate function she manages to reach.

We own, on average, over ten thousand things.\textsuperscript{25} Most of them are strangers to us within a year.

\bigskip
\begin{center}
\footnotesize
\renewcommand{\arraystretch}{1.15}
\begin{tabular}{@{}p{46mm}r@{}}
\toprule
objects per household & over 10\,000\textsuperscript{25} \\
plastic per year & 380\,000\,000 tonnes\textsuperscript{42} \\
uses per T-shirt & 7\textsuperscript{43} \\
breakdown time, cup & 450 years\textsuperscript{28} \\
CO\textsubscript{2} in the atmosphere & 427 ppm\textsuperscript{29} \\
topsoil lost per year & 24\,000\,000\,000 tonnes\textsuperscript{30} \\
planetary boundaries exceeded & 6 of 9\textsuperscript{44} \\
species threatened with extinction & 1\,000\,000\textsuperscript{45} \\
wild vertebrates since 1970 & down 69\,\%\textsuperscript{46} \\
forest lost per year & 10\,000\,000 hectares\textsuperscript{47} \\
humans in inadequate housing & 1\,800\,000\,000\textsuperscript{31} \\
humans without basic sanitation & 2\,400\,000\,000\textsuperscript{48} \\
humans driven from their homes & 120\,000\,000\textsuperscript{49} \\
humans vulnerable to climate change & 3\,500\,000\,000\textsuperscript{50} \\
active armed conflicts & 56\textsuperscript{51} \\
\bottomrule
\end{tabular}
\end{center}
\bigskip

We cannot say we did not know. Single-axis optimisation is not a moral betrayal; it is a structural property of narrow fields. And structural properties can be changed.

The design fields have no common grammar. No precise definition of what a form is, of what acts on it, of what care is. That has never been enough. But now we have the answers. The generative model in proposition 8 gives form an explicit expression: a probability distribution over the morphospace, weighted by each agent's estimate of the landscape, bounded by the grammar and the light cone. The formula does not say what is good. It says what is possible, what is probable, and where the blindness lies. The designer reading this has a choice. Not a choice between good and evil; a choice about how much of the world you wish to carry into what you make. Where you choose to see, you choose to attend. Where you choose to look away, you write it into the material. What you offer the world is never only the object. It is what the object does to everything it touches; the earth it lies in, the water it dissolves in, the body that bears it.

A cell in the body that bears it does not know you exist. It works, nonetheless, its whole life to hold you up; millions of divisions per second, towards a goal it will never see and never understand. It is a form-maker. What it builds is you, but the reference for what it builds does not lie within it; it lies in a system it has no empirical access to. It still builds rightly. Still billions of such cells rely on the whole being real and on the cooperation holding. They work out of care. That is not metaphorical; it is the only word the language gives us for a system that holds around itself because it needs to. So works every form-maker in the world. The bee does not know how the hive is. The root tip does not know where the trunk stands. The microorganism in the gut does not know that you are ill. They are cognitive glue in systems they cannot see.\textsuperscript{85,\,90} The interstitium, the connective tissue that for centuries was treated as passive barrier, was in 2018 reclassified as a fluid-filled organ with its own metabolism;\textsuperscript{107} a whole organ invisible because the language for it was missing. The invisible landscape became workable when the language was there. The craftsperson works this way. The farmer. The cook. Wherever a smaller system places form into a larger one it cannot oversee, it is the same structure, driven by the same ground: trust that the larger system is there and receives.

Every form-maker works for something larger than herself. The cell for the organ, the bee for the hive, the root tip for the tree, the craftsperson for the use that comes after. The structure is universal; it is not a human distinction. We do not have the apparatus to say what insight the cell or the tree has into the whole to which they belong. What we can say is that refusal is possible for us in a way it does not seem possible for them. The cell builds what it builds. The tree grows where it grows. The designer can choose not to be a designer, or to be one in name but not in deed. Responsibility comes from this asymmetry in choice, not from an asymmetry of insight. You are one of billions of such cooperations; without them, no you. This dependency is not a sentimental condition. It is the most fundamental reality in biology, and it repeats at every level upward. You are the cell for the organ, the organ is the cell for the body, the body is the cell for the family, for the city, for the species, for the biosphere. Each form-maker sits in a hierarchy of systems it holds up and that hold it up. The designer lives under the same structure as the cell; the difference is that she can step out of it.

All fields working with living systems know this. Ecology is built on the principle that everything hangs together. Systems dynamics is built on the principle that an intervention at one node propagates through the entire net. Medicine knows that the body cannot be cut into parts without loss. Climate science knows that the emission in one part of the atmosphere becomes my problem. The research ethics of these fields follows directly from the recognition: you may not do as you please, because your action does not stop where you stop seeing. Any foundation compatible with the present understanding of cognition, ecology, biology, and physics would inevitably produce an ethics that placed constraints on the designer. The design field has chosen not to know this. Not because the knowledge is not there; it is there in every neighbouring field and has been available for generations. But because a field that takes this knowledge in would have to stop doing what the field does today. We have escaped the recognition by not having a grammar that forces it on us.

The everyday form of the same pattern is the attention economy: service design at industrial scale, optimised for one axis. The harms are documented, large, and irreversible within a generation.\textsuperscript{94,\,95}

\begin{center}
\footnotesize
\renewcommand{\arraystretch}{1.15}
\begin{tabular}{@{}p{58mm}r@{}}
\toprule
average screen time, adolescents 13--18 & 7--9 h/day\textsuperscript{94} \\
rise in depression, young women USA 2010--2020 & +145\,\%\textsuperscript{94} \\
rise in self-harm admissions, girls 10--14 USA & +188\,\%\textsuperscript{94} \\
rise in anxiety among youth globally, PISA 2000--2018 & 36 of 37 countries\textsuperscript{94} \\
sleep under 7 hours per night, US adolescents & 31\,\%\,\msym{→}\,46\,\%\textsuperscript{94} \\
Meta's internal research, Instagram and teenage girls & 1 in 3 worse body image\textsuperscript{95} \\
internal research on recruitment of children under 13 & documented\textsuperscript{95} \\
children using TikTok per day globally & over 150\,000\,000 \\
Cambridge Analytica, profiled voters without consent & 87\,000\,000\textsuperscript{96} \\
time from front camera (2010) to mental-health collapse & 18 months\textsuperscript{94} \\
\bottomrule
\end{tabular}
\end{center}

The platform architecture makes Echo and Narcissus\textsuperscript{97} structural preconditions for participation: the nymph who can only repeat others, and he who drowns in his own reflection. Whoever gave form to that system gave form to that state. The numbers are not side effects. They are the primary effect of a design choice. Those who made the choice are called designers. They go to the same conferences. They win the same prizes. The blame is not individual. It is collective and structural, and therefore it can be cured only collectively and structurally.

A design theory that listens would acknowledge that systems hang together. It would acknowledge that each form sits at a place in a hierarchy, and that the place the form stands determines how far the consequences travel. A chair stands in a room, in a house, in a city, in an economy, in a biosphere. An algorithm stands in an app, in a phone, in a head, in a social circle, in a society. To operationalise this placement is not difficult. It is to ask, for each form: where in the system does it stand, what other layers does it touch, who works so that it can exist, who bears the cost. These questions have standard answers in ecology, in systems dynamics, in public health, in materials science. They need only to be made part of our craft. To treat the distinctions form and function, nature and artefact, biological and cultural, form-maker and user as gradients rather than as categories, as the foundation section sketches, is not a rhetorical move; it is an operational one. Research ethics then follows of itself. The designer then becomes one agent among many, rather than the sole ruler.

Practically this is a normalisation, not a revolution. Every deliverable must include a written justification of the axes that were represented and those that were chosen away, who bears the cost, which neighbouring fields were consulted. A medicine without efficacy data is not a medicine; it is a bottle. A design without such a justification is not a design; it is a proposal. Refuse commissions that cannot be justified in this way. Targeting systems for state violence. Attention extraction from children. Surveillance architectures against civilian populations. Interfaces made to produce dependency. An individual no loses commissions; a profession that says no gets to make a new one. That is the function of the organ: to turn the individual no into a professional no. Learn the substrates. Spend time with the mycelium, with the leather, with the code, with the cells that hold you up, with people who have lived in a system for twenty years. Competency never lies in the representation alone; it must be experienced. Extend the education. Not with a new course on ethics; the ethics follows from the rest. With a calibration of what the candidate has seen: what substrate she has worked in, what failure she has stood in, what neighbouring field she has consulted, what decisions she has made in which she said no. A candidate who has never said no to a commission has not completed her education. She has completed a copying.

These agents already exist. Mycelium packaging replaces polystyrene at concrete producers.\textsuperscript{52} Repair workshops are growing in Nordic cities. Slow computing, the free software movement, Ginkgo Bioworks, Ecovative, the craft tradition that never stopped; all of them already work under the same coordinate system, without a common name. They need not unite in one movement. They need only to be given a name for what they do.

\begin{center}
\includegraphics[width=0.85\linewidth]{larve.png}
\end{center}

Inside every caterpillar lie imaginal discs, small groups of undifferentiated cells that have lain inactive throughout the larval stage. When the larva enters the pupa, most of the body dissolves into a nutrient-rich mass. The discs activate, divide, and build an entirely new organism from the dissolved material. The brain is torn down and rebuilt. And still the butterfly remembers what the larva learned.\textsuperscript{27} What survives is the generalised insight, remapped to a new body. The metamorphosis is not a metaphor; it is a fact about what living substrate is capable of when given the opportunity to reorganise. It is the pattern fields follow through paradigm shifts;\textsuperscript{105} reorganisation of substrate, not merely argument. Design stands before the same transition. What survives the metamorphosis is not the exact operation. It is the care.

The cell does its work; the designer can choose whether she does hers. That is the whole difference, and it is the whole responsibility. A movement is a number of agents navigating the same gradient without central coordination: a common coordinate system, enough agents. The rest happens as mycelium through substrate. Form does not lie.

}
\end{widebody}
